% Cette fiche d'activité à été développée par
% + Erwan Tanguy-Legac <erwan.tanguy-legac@ens-rennes.fr>
% + Héloïse Troublé <heloise.trouble@ens-rennes.fr>

% Le template a été largement inspiré de celui de
% + Santiago Bautista <Santiago.Bautista@ens-rennes.fr>
% + Benjamin Bordais <Benjamin.Bordais@ens-rennes.fr>

% Elle est distribuée sous la licence Creative-Commons by-nc-sa 4.0
% (Attribution, Non-Commercial, Share Alike)
% qui peut être trouvée dans
% [le site de Creative Commons](https://creativecommons.org/licenses/by-nc-sa/4.0/)

\documentclass[12pt, a4paper]{article}
\usepackage[french]{babel}
\usepackage[utf8]{inputenc}
\usepackage[top=0.9in, left=0.6in, right=0.7in, bottom=1in]{geometry}
\usepackage{amssymb}
\usepackage{color}
\usepackage{fancyhdr}
\usepackage{longtable}
\usepackage{array}
\usepackage{titling}

\pagestyle{fancy}

\renewcommand{\headrulewidth}{1pt}
\chead{}
\lhead{Journal du Pirate - Plan de l'activité}
\rhead{Héloïse Troublé et Erwan Tanguy-Legac}

\renewcommand{\footrulewidth}{1pt}
\cfoot{\thepage} 
\lfoot{2024}
\rfoot{ENS Rennes}

\fancypagestyle{plain}{
	\fancyhf{}
	\cfoot{\thepage}
	\renewcommand{\headrulewidth}{0pt}
	\renewcommand{\footrulewidth}{1pt}}

\begin{document}
	
	\title{Journal du Pirate}
	\author{Héloïse Troublé et Erwan Tanguy-Legac}
	\date{12 avril 2024}
	\maketitle
	
	\paragraph*{Niveau :} $6^{\grave{e}me}$
	\paragraph*{Durée :} 55 minutes
	\paragraph*{Prérequis :} Aucun
	\paragraph*{Description du jeu :}
    L'activité consiste à modéliser des objets (îles) selon une description et à vérifier la présence d'une propriété (chemin).
	\paragraph*{Objectif pédagogique :}
    Introduction à la modélisation et à la vérification
  \paragraph*{Matériel nécéssaire :}
    Pour chaque groupe de 4 élèves, imprimer autant de fois qu'il y a d'îles (6 pour nous) de fiche des lieux.
	\begin{longtable}{c|m{3.9cm}|m{8.4cm}|m{2.4cm}}		\textbf{Durée} & \textbf{Phase} & \textbf{Activités et consignes} & \textbf{Organisation et matériel} \\ \hline\hline
		
		5' &
		
		Introduction
		
		&
		
		On se présente, on donne le contexte de l'activité.
		
		&
		
		A l'oral au tableau
		
		\\ \hline
		
	  5' &
		
    Modélisation du chemin
		
		&
	
    On leur demande de décrire le plus précisément possible le chemin décrit.
		
		& Par groupes de 4 (en îlots)
		
		\\ \hline
		
		5' &
		
		Mise en commun de la modélisation
		
		&
		
    On fait une mise en commun pour s'assurer qu'il n'y a aucun problème de compréhension du journal.
    On essaie de faire ressortir et comprendre les bons et mauvais points des différentes approches.

		&
		
		Au tableau
		
		\\ \hline
		
		10'
		
		&
		
	  Découverte des deux premières îles
		
		&
		
		On distribue les deux premières îles (île scorpion et île tortue). On essaie de passer dans les groupes et de voir le niveau de
    compréhension de chaque groupe pour savoir auxquels on va donner des extensions.
		
		& Par groupes de 4 (en îlots)
		
		\\ \hline
		
		5' &
		
	  Mise en commun
		
		&
		
    On met en commun pour que tous les élèves soient au même niveau.
    On leur demande d'expliquer exactement comment ils ont fait pour s'assurer que le trésor n'était pas présent sur les îles.
    On s'attend à des justifications avec des schémas.
		
		&
		
    Au tableau
		
		\\ \hline
		
		15'
		
		&
		
		Distribution des 4 autres îles

		&
		
    On distribue 3 ou 4 îles à chaque groupe (île de la fortune, île des oiseaux rapides,
    île des singes hurleurs, île des dauphins roses), en fonction de leur vitesse et compréhension lors des parties précédentes.
		
		&
		
		Par groupes de 4 (en îlots)

    \\ \hline

    10'

    &

    Institutionnalisation

    &

    On reprend l'attention des élèves, on fait l'Institutionnalisation : on explique l'intérêt de modéliser des propriétés et de les vérifier.
		
    &

    Au tableau

	\end{longtable}
	
\subsection*{Présentation :}
Des explorateurs ont trouvé le journal d’un grand pirate qui contient sans doute l’emplacement de son immense trésor. Vous devez les aider à
savoir sur quelle île il a été caché. Le journal contient uniquement la description du chemin qu’a emprunté le pirate. On va dans un premier
temps vous demander de schématiser le chemin, puis on vous distribuera des descriptions d’îles sur lesquelles le trésor pourrait être et on
vous demandera de trouver sur quelle île le trésor est.

\subsection*{Institutionnalisation :}
C’est de l’informatique, parce que grâce à nos cartes, on a réussi à modéliser toutes les îles et le chemin suivi par le pirate, pour vérifier
sur quelle île ce chemin était possible. Pour faire ça, on a choisi un format pour décrire les îles : nos cartes c’est ce qu’on appelle des
automates en informatique (ensemble d’états, ici les lieux (plage…), et des transitions entre des états, ici des chemins (chemin fleuri, scabreux…)).

\paragraph*{}
On a fait deux étapes de travail : \\
- D’abord modéliser les îles : choisir les bons états (les bons éléments sur chaque île), puis les bonnes transitions (les bons chemins entre
les éléments, avec parfois une information sur le temps qu’un chemin prend pour être parcouru) \\
- Puis une étape de vérification : pour chaque île, on a vérifié ce qu’on appelle des propriétés, c’est des caractéristiques qui permettent de
décrire les îles et voir si ça correspond au chemin du pirate, comme par exemple :
\\
“On peut débarquer sur cette île” (sinon c’est compliqué de l’explorer…), \\
“Il y a une forêt de palmier sur cette île”, \\
“La grotte est accessible depuis le volcan”, \\
“Le chemin jusqu’à la grotte prend 8h depuis la plage”... \\

Il n’y a qu’une seule île qui permettait de suivre le chemin du pirate dans le même temps que lui, on a donc pu déterminer sur quelle île
était caché le trésor.

En informatique, c’est un domaine qui s’appelle la vérification de modèle. Et ça sert par exemple au diagnostic : c’est exactement comme
pour les humains, on a envie de savoir si une machine fonctionne normalement. Pour ça, on représente les états dans lesquels elle peut être
(schéma débile au tableau en même temps), et comme on veut vérifier que la machine ne peut pas tomber en panne, on regarde si on peut aller
dans l’état “panne” en suivant les transitions au cours du fonctionnement de la machine. Si c’est pas le cas, on dit que l’état n’est pas
accessible, et on peut faire confiance à la machine !

Là on a fait tout ça à la main, mais on peut demander à un ordinateur de le faire : on automatise la vérification, mais au lieu de faire
confiance aveuglément à ce que l’ordinateur nous dit, on a un outil (les automates) pour pouvoir s’assurer que le résultat est bon. Ici,
on a pas eu besoin de vous croire sur parole pour savoir sur quelle île était le trésor, on est sûr grâce aux cartes !

\end{document}
